\begin{figure}[htbp]
\centering
\begin{tikzpicture}[>=Stealth, line width=0.8pt, font=\small,
nodo/.style={circle, fill=white, draw=black,
	inner sep=0pt, minimum size=5pt, line width=0.8pt},
cint/.style={circle, fill=white, draw=black,
	inner sep=0pt, minimum size=6pt, line width=0.8pt},
fantasma/.style={thin, gray, dashed},
tacca/.style={thin, gray!70}]

\def\a{1.0}
\def\tw{0.22}\def\thA{0.38}\def\thR{0.32}
\def\bw{0.16}\def\gw{0.30}\def\sr{0.08}

%% riferimento indeformato, tratteggiato
\draw[fantasma] (0,0) -- ({5*\a},0);

%% configurazione del meccanismo: C1 fermo, C2 ruotato attorno a C
\draw[line width=2.2pt] (0,0) -- ({3*\a},0);
\draw[line width=2.2pt] ({3*\a},0) -- ({5*\a},-0.6);

\node[above=2pt] at (0,0)      {$A$};
\node[above=2pt] at ({2*\a},0) {$B$};
\node[above=2pt] at ({3*\a},0) {$C$};
\node[below right=1pt] at ({5*\a},-0.6) {$D$};

\node at ({1.4*\a},0.28)  {$\mathcal{C}_1$};
\node at ({4.3*\a},-0.2)   {$\mathcal{C}_2$};

%% angolo virtuale theta_2 in C
\draw[tacca] ({3.2*\a-0.35},0) -- ({3.2*\a+0.35},0);
\draw[thin] ({3.2*\a+0.3},0) arc (0:-16.7:0.3);
\node[right=1pt] at ({3*\a+0.32},-0.4) {$\vartheta_2$};

%% cerniera in A e carrello in B, invariati
\draw (0,0) -- ++(-\tw,-\thA) -- ++(2*\tw,0) -- cycle;
\draw (-\gw,-\thA) -- (\gw,-\thA);
\foreach \x in {-0.24,-0.12,0,0.12,0.24}
\draw[thin] (\x,-\thA) -- ++(-0.09,-0.11);
\node[nodo] at (0,0) {};

\draw ({2*\a},0) -- ++(-\tw,-\thR) -- ++(2*\tw,0) -- cycle;
\draw ({2*\a-\gw},-\thR) -- ({2*\a+\gw},-\thR);
\draw ({2*\a-\bw},-\thR-0.12) circle (2.4pt);
\draw ({2*\a+\bw},-\thR-0.12) circle (2.4pt);
\draw ({2*\a-\gw},-\thR-0.24) -- ({2*\a+\gw},-\thR-0.24);
\foreach \x in {-0.24,-0.12,0,0.12,0.24}
\draw[thin] ({2*\a+\x},-\thR-0.24) -- ++(-0.09,-0.11);
\fill[white] ({2*\a-\sr},0) arc(180:360:\sr) -- cycle;
\draw ({2*\a-\sr},0) arc(180:360:\sr);

\node[cint] at ({3*\a},0) {};

%% carrello in D soppresso: fantasma nella posizione originaria
\draw[fantasma] ({5*\a},0) -- ++(-\tw,-\thR) -- ++(2*\tw,0) -- cycle;
\node[gray, above=2pt] at ({5*\a},0) {\footnotesize(soppresso)};
\node[nodo] at ({5*\a},-0.6) {};

%% reazione incognita Y_D, verso l'alto, applicata nel punto attuale
\draw[-{Latex[length=2.8mm]}, thick] ({5*\a},-1.2) -- ({5*\a},-0.65);
\node[below=2pt] at ({5*\a},-0.95) {$Y_D$};

%% carico F, verso il basso, applicato nel punto attuale
\draw[-{Latex[length=2.8mm]}, thick] ({4*\a},0.7) -- ({4*\a},0.1);
\node[right=2pt] at ({3.4*\a},0.5) {$F$};

\end{tikzpicture}
\caption{Il meccanismo usato per determinare $Y_D$: si sopprime
(concettualmente) il carrello in $D$, in grigio nella sua posizione
originaria. Cerniera in $A$ e carrello in $B$ bloccano
$\mathcal{C}_1$, che resta fermo; $\mathcal{C}_2$ è allora libero di
ruotare, di un angolo virtuale $\vartheta_2$, attorno al punto fisso
$C$. Solo $F$ e la reazione incognita $Y_D$ compiono lavoro in
questo moto.}
\label{fig:gerber_mecc_D}
\end{figure}
